\chapter{The prediction problem and the system under investigation}
\label{ch:context}

\begin{saybox}[title=What to say at the start of this section]
Open by stating the problem in one sentence --- predicting how many days a
cheese formulation remains acceptable under given storage conditions ---
then explain that the system is not one model but six, and that this
structure matters for everything that follows, because it determines how
little data each individual model actually sees.
\end{saybox}

\section{The prediction task}

The system predicts \emph{shelf life in days} for a cheese formulation
under specified storage conditions. The target variable,
\texttt{shelf\_life\_days}, is continuous and strictly positive, so the
problem is posed as regression. Given a feature vector $\mathbf{x}$
describing the cheese matrix, the treatment applied to it, the storage
environment, and the spoilage or safety indicator being tracked, the model
estimates
\begin{equation}
\hat{y} = f(\mathbf{x}), \qquad
\hat{y} \in \mathbb{R}_{>0} \text{ (days)} .
\end{equation}

Two distinct questions are asked of the same physical system, and the
project treats them as two separate prediction tasks:

\begin{description}[leftmargin=1.6em, style=nextline, font=\bfseries\color{dgreen}]
\item[General shelf life] the time until a quality indicator crosses an
  acceptability threshold --- yeasts and moulds, total viable count,
  sensory acceptability, pH drift, peroxide value, oxidation index, visible
  moulds, or clostridial late blowing. This is a \emph{quality} endpoint.
\item[Safety endpoint] the time until a pathogen-related criterion is
  reached --- \emph{Listeria monocytogenes}, \emph{Staphylococcus aureus},
  \emph{Salmonella}. This is a \emph{safety} endpoint and, as
  \Cref{sec:censoring} discusses, its target semantics are not identical
  to those of the quality endpoint.
\end{description}

\section{Six specialists, not one model}

The production architecture routes each prediction to one of six
specialist models, indexed by the Cartesian product of cheese category and
prediction task:
\begin{equation}
\mathcal{S} = \{\text{soft}, \text{semi-hard}, \text{hard}\}
\times
\{\text{general\_shelf\_life}, \text{safety\_endpoint}\},
\qquad |\mathcal{S}| = 6 .
\end{equation}

This was verified against the implementation, not assumed: the training
driver \evd{train\_specialists.py} iterates over exactly these three
categories and two tasks, loading one category file per category and
filtering it on the \texttt{model\_task} column. Each specialist is trained
independently, with its own feature schema detection, its own preprocessing
fit, and its own set of four candidate model families.

\fig{fig04_dataset_composition}{0.92}{%
  Composition of the V7 training corpus. The three category files contain
  34\,000 rows in total, partitioned into six specialist training sets.
  Every row carries \texttt{data\_origin = synthetic\_literature\_%
  constrained\_v7}; there is no real-observation subset in this corpus.
  Generated from the three raw specialist CSV files.}

The practical consequence of this architecture is that the effective
training set for any one model is much smaller than the headline corpus
size. The hard/safety-endpoint specialist, for example, trains on
$2\,240$ rows, not on $34\,000$. This becomes important in
\Cref{ch:synthetic}, where the motivation for synthetic expansion is
examined quantitatively.

\section{The four candidate model families}

For each of the six specialists, four model families are trained and
compared:

\begin{description}[leftmargin=1.6em, style=nextline, font=\bfseries\color{dgreen}]
\item[Random Forest] a bagged ensemble of unpruned regression trees;
  variance reduction through bootstrap resampling and feature subsampling.
\item[LightGBM] gradient boosting with leaf-wise (best-first) tree growth
  and histogram-binned split finding.
\item[XGBoost] gradient boosting with depth-wise tree growth and explicit
  $L_1$/$L_2$ leaf-weight regularization in the objective.
\item[Explainable Boosting Machine (EBM)] a glass-box generalized additive
  model with a bounded number of pairwise interaction terms, cyclically
  boosted one feature at a time.
\end{description}

Their exact production hyperparameters are recovered from source and
reported in \Cref{tab:hparams}; the mathematics of each is developed in
\Cref{ch:rf,ch:lgbm-xgb,ch:ebm}.

\begin{table}[htbp]\centering\small
\caption{Production hyperparameters, read directly from
  \texttt{train\_specialists.py}. The seed is global
  (\texttt{SEED = 42}) and shared by the splitter and every model.}
\label{tab:hparams}
\begin{tabular}{lp{9.4cm}}
\toprule
Model & Configuration as implemented \\
\midrule
Random Forest &
\texttt{n\_estimators=400}, \texttt{max\_depth=None},
\texttt{min\_samples\_leaf=2}, \texttt{random\_state=42}, \texttt{n\_jobs=-1} \\
\addlinespace[2pt]
LightGBM &
\texttt{n\_estimators=2000}, \texttt{learning\_rate=0.03},
\texttt{num\_leaves=31}, \texttt{min\_child\_samples=20},
\texttt{subsample=0.9}, \texttt{colsample\_bytree=0.9};
early stopping after 50 rounds on the validation set \\
\addlinespace[2pt]
XGBoost &
\texttt{n\_estimators=2000}, \texttt{learning\_rate=0.03},
\texttt{max\_depth=6}, \texttt{subsample=0.9},
\texttt{colsample\_bytree=0.9}, \texttt{eval\_metric="rmse"};
early stopping after 50 rounds \\
\addlinespace[2pt]
EBM &
\texttt{interactions=10}, \texttt{max\_bins=256},
\texttt{learning\_rate=0.02}, \texttt{max\_rounds=3000},
\texttt{min\_samples\_leaf=4}, \texttt{outer\_bags=8} \\
\bottomrule
\end{tabular}
\end{table}

\begin{findingbox}[title=Note on capacity]
Two of the four production models use \emph{two thousand} boosting rounds
with early stopping, and the Random Forest grows 400 fully-developed trees
with a minimum leaf size of two. These are high-capacity configurations.
If the high scores were produced by capacity, reducing capacity should
reduce them. \Cref{ch:diagnostics} tests exactly that, and it does not.
\end{findingbox}

\section{Feature space}

Schema detection is performed per specialist at training time by
\texttt{detect\_schema()}: any column that is not the target and not on the
exclusion list is classified as numeric, binary, or categorical by
inspecting its dtype and its set of unique values. The V7 files carry 63
columns, of which the identifier, provenance, and zero-variance columns are
excluded, leaving the physically meaningful inputs.

The features that reach the model describe four groups:

\begin{itemize}[leftmargin=1.3em, itemsep=1pt]
\item \textbf{Matrix composition:} \texttt{matrix\_ph},
  \texttt{matrix\_water\_activity}, \texttt{matrix\_moisture\_pct},
  \texttt{matrix\_fat\_pct}, \texttt{matrix\_protein\_pct},
  \texttt{matrix\_salt\_pct}, \texttt{matrix\_ripening\_days},
  \texttt{pasteurization\_applied}, \texttt{food\_matrix},
  \texttt{physical\_form}. (\texttt{base\_cheese\_name} is \emph{excluded}
  as redundant with \texttt{food\_matrix}; it is retained in the data and
  used throughout this report to label results.)
\item \textbf{Storage environment:} \texttt{storage\_temperature\_c},
  \texttt{packaging\_type}, \texttt{headspace\_oxygen\_pct},
  \texttt{headspace\_co2\_pct}, \texttt{headspace\_n2\_pct}.
\item \textbf{Treatment:} \texttt{treatment\_type},
  \texttt{application\_method}, \texttt{ingredient\_count},
  \texttt{primary\_ingredient\_name}, \texttt{primary\_ingredient\_family},
  \texttt{canonical\_concentration\_value}, \texttt{is\_control}. (The raw
  \texttt{primary\_concentration} and its unit column are excluded in
  favour of the unit-harmonised canonical value.)
\item \textbf{Indicator definition:} \texttt{indicator\_group},
  \texttt{indicator\_type}, \texttt{indicator\_threshold},
  \texttt{initial\_indicator\_value}, and, for safety specialists only,
  \texttt{initial\_inoculum\_log\_cfu\_g} and \texttt{growth\_support}.
\end{itemize}

\subsection{What is deliberately excluded, and why it matters}
\label{sec:exclusions}

The exclusion list is not a formality. Reading it is one of the more
informative things one can do in this codebase, because the comments record
checks that were actually performed. Eleven V7-specific columns are
excluded, each for a stated and verifiable reason: \texttt{routing\_category}
is $1$:$1$ with \texttt{cheese\_category};
\texttt{physical\_form\_observed} is constant zero;
\texttt{endpoint\_threshold\_basis} and \texttt{is\_challenge\_test} are
bijective with \texttt{model\_task} and therefore constant within any one
specialist; \texttt{evidence\_class}, \texttt{observed\_fields},
\texttt{source\_url}, \texttt{source\_note},
\texttt{generation\_rule\_version}, and \texttt{matrix\_profile\_confidence}
are provenance rather than physics.

Critically for this investigation, \texttt{source\_rule\_id} itself is
excluded from the feature set. The models never see the rule label. The
problem identified in \Cref{ch:discovery} is therefore \emph{not} that the
rule identifier leaks into the features --- it does not --- but that the
\emph{evaluation split} ignores it.

\begin{limitbox}[title=One exclusion worth discussing with the panel]
\texttt{target\_is\_lower\_bound} is excluded with the comment that it
``describes the \textsc{target}'s censoring semantics, not a live input a
user could supply.'' That reasoning is sound for a deployed predictor:
the flag would not be known at prediction time. But it means the
regression objective treats censored and uncensored targets identically
during fitting. \Cref{sec:censoring} returns to this.
\end{limitbox}
